\section{The \method{} Framework with \sgpn}
\label{sec:method}

% ─────────────────────────────────────────────────────────────────────────────
\subsection{Theoretical Foundations}
\label{sec:theory}

We establish four results that jointly guarantee both the
\emph{soundness} (every generated circuit is valid) and
\emph{completeness} (every valid circuit is reachable) of our
safety-masked sequential generation framework.
Worked $2$-variable examples illustrating each result, and all proofs, are
given in the supplementary material (App.~A--B).

\begin{theorem}[Monotonicity of Conducting Paths]
\label{thm:mono}
For any transistor network $G$ and any transistor $e \notin E(G)$,
$\mathrm{Paths}(G) \subseteq \mathrm{Paths}(G \cup \{e\})$.
\end{theorem}
\begin{corollary}[Incremental Safety Checking]
\label{cor:safety}
If $G$ is safe, then $G \cup \{(u,v,\ell)\}$ is safe iff $(u,v,\ell)$ alone
does not create a new $\SRC \to \SNK$ path under any $\pi \in \Poff$.
\end{corollary}
This reduces safety verification per candidate from $O(|E|)$ to
$O(|\Poff|\cdot|V|)$, independent of circuit size.

\begin{theorem}[SOP Network Correctness]
\label{thm:sop}
Let $G_{\mathrm{SOP}}(f)$ be the series-parallel network with one
$N$-transistor series chain per $\pi \in \Pon$, all sharing $\SRC$/$\SNK$.
Then $G_{\mathrm{SOP}}(f)$ correctly implements $f$.
\end{theorem}
\begin{theorem}[Completeness of Safety-Masked Search]
\label{thm:complete}
For any valid network $G^*$ implementing $f$, every permutation of
$E(G^*)$ is a valid trajectory under the safety mask.
Equivalently, the safety-masked MDP reaches $G^*$ from the empty graph
in $|E(G^*)|$ steps following any edge ordering.
\end{theorem}
\paragraph{Implication for search completeness.}
Theorem~\ref{thm:complete} is the key correctness guarantee of our framework:
\emph{the safety mask never prunes any valid solution}.
It only forbids transistors whose addition immediately creates an off-pattern
path — those that cannot appear in any correct, safe network.
The policy therefore searches the full space of valid partial sub-networks;
no optimal solution is made unreachable by the masking procedure.

\subsection{MDP Formulation}
\label{sec:mdp}

\paragraph{State.}
The state at step $t$ is $\mathcal{G}_t = G_t \cup C$, where $G_t$ is
the partial synthesized network and $C = G_{\mathrm{SOP}}(f)$ is the
fixed compensate network from Theorem~\ref{thm:sop}.
The two share only $\SRC$ and $\SNK$.

\paragraph{Action.}
Action $a_t = (u,v,\ell)$ adds one transistor.
Two masks apply: (i) \textbf{reachability mask} on $u$ prunes structurally
disconnected nodes; (ii) \textbf{safety mask} on $(u,v,\ell)$ forbids
transistors creating $\Poff$-paths (Corollary~\ref{cor:safety}).
By Theorem~\ref{thm:complete}, the safety mask is complete.

\paragraph{Reward.}
Terminal reward only (no per-step penalty):
\begin{equation}
  r = \begin{cases}
    T^*(f)/T   & \text{Phase 2: complete network, } T \text{ transistors}\\
    \Phi^*(f)/\Phi & \text{Phase 3: complete PDN+PUN pair}\\
    0          & \text{incomplete.}
  \end{cases}
  \label{eq:reward}
\end{equation}
$T^*(f)$ and $\Phi^*(f)$ are self-tightening best values found so far for $f$.

% ─────────────────────────────────────────────────────────────────────────────
\subsection{Unified Graph State Representation}

At each step $t$, the policy observes the unified graph
$\mathcal{G}_t = G_t \cup C$, where:
\begin{itemize}[leftmargin=*, topsep=1pt, itemsep=0pt]
  \item $G_t = (V_t, E_t)$ is the \emph{partial synthesized network}
        built so far (edges are placed transistors with literal labels).
  \item $C = G_{\mathrm{SOP}}(f)$ is the \emph{fixed compensate network}
        (Theorem~\ref{thm:sop}), kept constant across all steps of the episode.
\end{itemize}
The two components share only $\SRC$ and $\SNK$; all internal nodes
are disjoint.
This unified representation gives the GNN simultaneous access to
the current circuit topology ($G_t$) and the full logical specification
of $f$ ($C$), enabling the policy to reason jointly about what has
been built and what coverage obligations remain.

\paragraph{Node and edge features.}
Each node $v \in \mathcal{G}_t$ carries a feature vector encoding:
node type ($\SRC$, $\SNK$, $G$-internal, $C$-internal),
connectivity degree in $G_t$ and in $C$ separately,
and a coverage indicator (whether $v$ lies on a chain for an
already-covered on-pattern).
Each edge $(u, v, \ell) \in E_t \cup E_C$ carries a feature vector
encoding: component membership ($G$ vs.\ $C$), variable index,
and polarity.
Node feature dimension $d_\mathrm{node}$ and edge feature
dimension $d_\mathrm{edge}$ are given in the supplementary (App.~C).

\paragraph{Virtual new-node embedding.}
The node space is extended to $V_t^+ = V_t \cup \{v_\mathrm{new}\}$
by appending a single learnable embedding $\mathbf{e}_\mathrm{new} \in \mathbb{R}^d$
for the virtual new-node option.
This allows the policy to propose a fresh internal node at any step
without a separate node-count decision, keeping the action space
consistent regardless of how many nodes have been created so far.

\subsection{\sgpn: Five-Head Factored Policy}

\begin{figure}[t]
\centering
% ── Figure: SGPN Architecture ─────────────────────────────────────────────
\begin{tikzpicture}[
    >=Stealth,
    font=\small,
    bigbox/.style = {draw, rounded corners=4pt, minimum width=7.6cm,
                     minimum height=0.65cm, align=center},
    head/.style   = {draw, rounded corners=3pt, minimum width=1.6cm,
                     minimum height=0.85cm, align=center, fill=blue!8,
                     font=\scriptsize},
    mask/.style   = {draw, rounded corners=2pt, minimum width=1.55cm,
                     minimum height=0.36cm, align=center, fill=orange!20,
                     font=\scriptsize},
    arr/.style    = {->, thick},
    lbl/.style    = {font=\scriptsize, text=gray},
]

%% ── Input ─────────────────────────────────────────────────────────────────
\node[bigbox, fill=green!8] (ug) at (0, 0)
  {Unified Graph $\;H_t = G_t \;\cup\; C_g$};
\node[lbl, left=0.15cm of ug] {input};

%% ── MPNN ──────────────────────────────────────────────────────────────────
\node[bigbox, fill=gray!10, below=0.42cm of ug] (mpnn)
  {\textbf{MPNN} Backbone ($L$ layers)
   $\quad\mathbf{h}_i\!\in\!\mathbb{R}^d,\;$
   $\mathbf{g}\!\in\!\mathbb{R}^{2d}$};
\draw[arr] (ug) -- (mpnn);
\node[lbl, left=0.15cm of mpnn] {encode};

%% ── Four heads ────────────────────────────────────────────────────────────
\node[head, below=0.8cm of mpnn, xshift=-2.85cm] (h1)
  {Head~1\\Node-1};
\node[head, below=0.8cm of mpnn, xshift=-0.95cm] (h2)
  {Head~2\\Node-2};
\node[head, below=0.8cm of mpnn, xshift= 0.95cm] (h3)
  {Head~3\\Variable};
\node[head, below=0.8cm of mpnn, xshift= 2.85cm] (h4)
  {Head~4\\Sign};

\draw[arr] (mpnn.south) -- ++(0,-0.10) -| (h1.north);
\draw[arr] (mpnn.south) -- ++(0,-0.10) -| (h2.north);
\draw[arr] (mpnn.south) -- ++(0,-0.10) -| (h3.north);
\draw[arr] (mpnn.south) -- ++(0,-0.10) -| (h4.north);
\node[lbl, left=0.15cm of h1] {decode};

%% ── Safety masks ──────────────────────────────────────────────────────────
\node[mask, below=0.22cm of h1] (m1) {\textit{safety mask}};
\node[mask, below=0.22cm of h2] (m2) {\textit{safety mask}};
\node[mask, below=0.22cm of h3] (m3) {\textit{safety mask}};
\node[mask, below=0.22cm of h4] (m4) {\textit{safety mask}};
\draw[arr] (h1)--(m1); \draw[arr] (h2)--(m2);
\draw[arr] (h3)--(m3); \draw[arr] (h4)--(m4);

%% ── Action output ─────────────────────────────────────────────────────────
\node[bigbox, fill=red!7, below=0.5cm of m2, xshift=0.95cm] (act)
  {Action $(u^*,\, v^*,\, k^*,\, s^*)$: one transistor placement};
\draw[arr] (m1.south) -- ++(0,-0.15) -| (act.north);
\draw[arr] (m2.south) -- ++(0,-0.15) -| (act.north);
\draw[arr] (m3.south) -- ++(0,-0.15) -| (act.north);
\draw[arr] (m4.south) -- ++(0,-0.15) -| (act.north);
\node[lbl, left=0.15cm of act] {output};

\end{tikzpicture}

\caption{Architecture of the \sgpn (\SGPN).
The unified graph $\mathcal{G}_t = G_t \cup C_{\mathrm{SOP}}$ is encoded
by a shared MPNN backbone; four factored heads with per-head safety masks
decode the transistor-placement action $(u^*,v^*,k^*,s^*)$.}
\label{fig:sgpn}
\end{figure}

The \SGPN uses an MPNN backbone~\cite{gilmer2017mpnn} with $L$ message-passing
layers and hidden dimension $d$ to compute node and graph-level features from
$\mathcal{G}_t$.
Let $\mathbf{h}_i \in \mathbb{R}^d$ be the node feature for node $i$
and $\mathbf{g} \in \mathbb{R}^{2d}$ the graph-level summary.
Four specialized MLP heads decode the transistor-placement action from these
features, each operating over a different sub-decision of the factored action
space.
All four heads share the same MPNN forward pass; safety masks are applied
independently within each head.

\paragraph{Head~1: Node-1 (source node).}
Scores each $G$-node $u \in V_t^+$ as the source of the new transistor:
\[
  s_u = \mathrm{MLP}_1([\mathbf{h}_u \| \mathbf{g}]).
\]
Logits are masked to allow only nodes in the reachability set.

\paragraph{Head~2: Node-2 (drain node).}
Conditioned on the sampled source $u^*$, scores each node $v \in V_t^+$,
$v \neq u^*$:
\[
  s_v = \mathrm{MLP}_2([\mathbf{h}_{u^*} \| \mathbf{h}_v \| \mathbf{g}]).
\]
Logits are masked to exclude nodes that have no safe literal for any
$(u^*, v, \ell)$ (precomputed BFS safety check).

\paragraph{Head~3: Variable index.}
Given the sampled edge $(u^*, v^*)$, selects the control variable:
\[
  \mathbf{s}_{\mathrm{var}} = \mathrm{MLP}_3([\mathbf{h}_{u^*} \| \mathbf{h}_{v^*}])
  \in \mathbb{R}^{N_{\mathrm{var}}}.
\]
Logits are safety-masked: only variables with at least one safe polarity
for $(u^*, v^*)$ are valid.

\paragraph{Head~4: Polarity (sign).}
Selects positive or negative literal for the chosen variable:
\[
  \mathbf{s}_{\mathrm{sgn}} = \mathrm{MLP}_4([\mathbf{h}_{u^*} \| \mathbf{h}_{v^*}])
  \in \mathbb{R}^{2}.
\]
Logits are safety-masked per the chosen variable.

\paragraph{Factored log-probability.}
The log-probability of placing transistor $(u, v, k, s)$ factorizes as:
\begin{equation}
  \log \pi(a \mid \mathcal{G}_t)
  = \underbrace{\log p_1(u)}_{\text{Head 1}}
  + \underbrace{\log p_2(v \mid u)}_{\text{Head 2}}
  + \underbrace{\log p_3(k \mid u,v)}_{\text{Head 3}}
  + \underbrace{\log p_4(s \mid u,v,k)}_{\text{Head 4}},
  \label{eq:logp}
\end{equation}
where all probabilities are masked softmax over safe candidates.
By Corollary~\ref{cor:safety}, every action sampled from this distribution
maintains the safety invariant; no post-hoc filtering is required.

\paragraph{Batched forward pass.}
During PPO training, $K$ trajectories are collected per function per update
step, each of length up to $T_{\max}$.
Naively, this requires $O(K \cdot T_{\max})$ separate MPNN forward passes.
Instead, all step graphs across all trajectories and all functions in the
mini-batch are packed into a single \texttt{PackedGraph} and processed with
\emph{one} MPNN call.
Per-step node features are then sliced from the shared output tensor for
the head MLPs.
This reduces the dominant computation cost by $K \cdot T_{\max}\times$
and makes large-batch PPO training feasible on a single GPU.

\subsection{Phase 1: NLL Pretraining}

We collect a dataset $\mathcal{D}$ of known-optimal transistor networks,
each paired with its target Boolean function.
For each network, we replay the transistor-placement sequence and record
the $(u, v, k, s)$ labels at every step.

\paragraph{Construction order (supervision linearization).}
\label{par:order}
A transistor network is an \emph{unordered} set of devices, so replaying it
as a sequence requires choosing a linearization
$e_1,\dots,e_{|E^*|}$, and that choice defines the learning problem.
A natural default is to order devices by node index (equivalently, by a
breadth-first traversal from the rails), but this makes the source label an
artifact of the indexing convention rather than of the circuit.
We instead use a \emph{path (ear) decomposition}: the sequence is emitted as
successive walks, where the first walk starts at \textsc{src} and extends
through unvisited nodes until it closes into already-visited structure
(preferentially at \textsc{snk}), and each later walk (an \emph{ear}) starts
at an already-visited node and extends until it reconnects.
Every emitted action therefore satisfies the invariant that $u$ is already
present in the partial network, so the walk direction fixes the orientation.

This ordering aligns supervision with the semantics of the task: because
conduction is achieved only when a \textsc{src}$\to$\textsc{snk} path
\emph{closes}, decomposing the target into completed paths makes each label
continue a coherent sub-goal.  It also changes the difficulty profile.
During a walk extension the partial network has exactly one dangling
endpoint, and that endpoint \emph{is} the next source, so Head~1 becomes
structurally determined rather than conventional; the residual difficulty
concentrates at \emph{ear starts}, where the model must choose which
existing node to branch from --- precisely the structure-sharing decision
that governs compactness.  Section~\ref{sec:exp} shows this single change
raises source-selection accuracy by roughly $30$ points, far more than any
architectural variation we tested.

\paragraph{Label consistency.}
A Boolean function typically admits many distinct optimal or near-optimal
networks.  Retaining several implementations of the \emph{same} function
pairs identical states with conflicting targets, giving the objective an
irreducible floor.  We therefore retain a single implementation per
function (the most compact one), which both removes the conflict and makes
the supervision consistent with the downstream objective of minimizing
transistor count.

The pretraining objective is the mean negative log-likelihood:
\[
  \mathcal{L}_{\mathrm{NLL}}
  = -\frac{1}{|\mathcal{D}|} \sum_{(G^*, f) \in \mathcal{D}}
    \sum_{t=1}^{|E^*|} \log \pi(a_t^* \mid \mathcal{G}_{t-1}),
\]
where $a_t^*$ is the ground-truth transistor at step $t$.
The unified graph $\mathcal{G}_{t-1}$ is constructed from the partial
prefix $G_{t-1}^*$ and the fixed compensate network $C_f$.
Pretraining is batched via \texttt{PackedGraph} (all samples in a mini-batch
share a single MPNN forward pass) and trained with Adam until convergence.

\subsection{Phase 2: PPO Finetuning}
\label{sec:rl}

After pretraining, we finetune with off-policy PPO.
A frozen copy of the current model (\emph{agent}) generates trajectories;
the live model recomputes log-probabilities for the PPO surrogate loss.

\paragraph{Reward design.}
We use a purely terminal reward with no per-step penalty.
Let $T$ be the number of transistors placed and $T^*$ the best transistor
count found so far for this function (\emph{adaptive self-tightening target}):
\[
  r = \begin{cases} T^*/T & \text{if complete,} \\ 0 & \text{otherwise.} \end{cases}
\]
Once the policy discovers a shorter circuit, $T^*$ updates automatically,
so subsequent solutions that are longer receive a reduced reward.
Each step in the trajectory receives the same scalar $r$ (no discounting).

\paragraph{GRPO baseline.}
For each function in the batch, we sample $K$ trajectories in parallel.
The advantage for each trajectory is normalized within the group:
\[
  \hat{A}_k = \frac{r_k - \bar{r}}{\sigma_r + \varepsilon},
  \qquad
  \bar{r} = \frac{1}{K}\sum_{k=1}^K r_k,\quad
  \sigma_r = \sqrt{\frac{1}{K}\sum_{k=1}^K (r_k - \bar{r})^2}.
\]
This group-relative policy optimization (GRPO)~\cite{grpo} baseline
eliminates EMA lag and gives a strong positive signal to any trajectory
shorter than its group average without requiring a value network.

\paragraph{Entropy bonus.}
To prevent mode collapse (where all $K$ trajectories converge to the same
solution), we add an entropy regularization term to the PPO loss:
\[
  \mathcal{L} = \mathcal{L}_{\mathrm{PPO}} - \lambda_H \cdot \bar{H},
\]
where $\bar{H}$ is the mean per-step policy entropy across all four heads
and $\lambda_H$ is a tunable coefficient (default 0.01).
The entropy bonus maintains exploration, ensuring the policy continues to
occasionally sample shorter-than-current-best circuits.

\paragraph{Batched MPNN forward.}
To reduce wall-clock time, all trajectory steps across all $K$ trajectories
and all functions in the batch are packed into a single \texttt{PackedGraph}
and processed with one MPNN forward pass.
Subsequent head MLPs slice per-step node features from the shared output.
This replaces $O(K \cdot T_{\max})$ sequential forward passes with~1.

\paragraph{Agent synchronization.}
The frozen agent is re-synchronized to the live model every $M$ PPO steps
to keep the importance-sampling ratio $\pi/\pi_{\mathrm{old}}$ within a
controlled range.

\begin{figure*}[t]
\centering
% ── Figure: Three-Phase Training Flow ─────────────────────────────────────
\begin{tikzpicture}[
    >=Stealth,
    font=\small,
    box/.style    = {draw, rounded corners=3pt, align=center},
    phase/.style  = {box, minimum width=4.2cm, minimum height=2.2cm},
    rwd/.style    = {box, minimum width=3.4cm, minimum height=0.45cm,
                     fill=yellow!25, font=\scriptsize},
    ckpt/.style   = {box, minimum width=2.2cm, minimum height=0.4cm,
                     fill=white, font=\scriptsize},
    arr/.style    = {->, thick},
    darr/.style   = {->, thin, dashed, gray},
    lbl/.style    = {font=\scriptsize},
]

%% ── Dataset ───────────────────────────────────────────────────────────────
\node[box, fill=gray!12, minimum width=1.6cm, minimum height=0.9cm,
      font=\scriptsize] (ds) at (-7.0, 0)
  {Dataset\\$\mathcal{D}$};

%% ── Phase 1 ──────────────────────────────────────────────────────────────
\node[phase, fill=blue!7]  (p1) at (-3.8, 0) {};
\node[font=\scriptsize\bfseries, above=0.5cm of p1.center] {\textbf{Phase 1}};
\node[font=\scriptsize, at=(p1.center), yshift= 0.1cm] {NLL Pretraining};
\node[rwd] (r1) at (-3.8, -0.72) {$\mathcal{L}_{\mathrm{NLL}}$};

\draw[arr] (ds.east) -- (p1.west);

\node[ckpt] (ck1) at (-3.8, -1.72) {$\pi^{(1)}_\theta$};
\draw[arr] (p1.south) -- (ck1.north);

%% ── Phase 2 ──────────────────────────────────────────────────────────────
\node[phase, fill=green!7] (p2) at (1.0,  0) {};
\node[font=\scriptsize\bfseries, above=0.5cm of p2.center] {\textbf{Phase 2}};
\node[font=\scriptsize, at=(p2.center), yshift= 0.28cm] {PPO Count RL};
\node[rwd] (r2) at (1.0, -0.35) {$r = T^*(f)\;/\;T$};
\node[font=\scriptsize, text=gray] at (1.0,-0.85) {GRPO + entropy};

\draw[arr] (ck1.east) -- ++(0.2,0) |- (p2.west);

\node[ckpt] (ck2) at (1.0, -1.72) {$\pi^{(2)}_\theta$};
\draw[arr] (p2.south) -- (ck2.north);

%% ── Phase 3 ──────────────────────────────────────────────────────────────
\node[phase, fill=red!6]  (p3) at (6.0,  0) {};
\node[font=\scriptsize\bfseries, above=0.5cm of p3.center] {\textbf{Phase 3}};
\node[font=\scriptsize, at=(p3.center), yshift= 0.28cm] {PPO Physical RL};
\node[rwd] (r3) at (6.0, -0.35) {$r = \Phi^*(f)\;/\;\Phi$};
\node[font=\scriptsize, text=gray] at (6.0,-0.85) {joint PDN\,+\,PUN};

\draw[arr] (ck2.east) -- ++(0.2,0) |- (p3.west);

\node[ckpt] (ck3) at (6.0, -1.72) {$\pi^{(3)}_\theta$};
\draw[arr] (p3.south) -- (ck3.north);

%% ── ASTRAN oracle ─────────────────────────────────────────────────────────
\node[box, fill=orange!15, minimum width=1.8cm, minimum height=0.7cm,
      font=\scriptsize] (astran) at (8.4, -0.7)
  {ASTRAN\\$\Phi$};

\draw[darr] (astran.west) -- node[above, font=\tiny]{$\Phi$} (p3.east);

%% ── Phase header labels ───────────────────────────────────────────────────
\node[lbl, text=blue!70,          above=1.32cm of p1] {initialise};
\node[lbl, text=green!50!black,   above=1.32cm of p2] {min.\ transistors};
\node[lbl, text=red!60,           above=1.32cm of p3] {optimise placement};

\end{tikzpicture}

\caption{Three-phase training pipeline of \method.
Phase~1 initialises the \SGPN by NLL imitation on known-optimal circuits.
Phase~2 finetunes toward minimum transistor count via PPO with adaptive
reward $r = T^*(f)/T$ and GRPO baseline.
Phase~3 jointly synthesises PDN and PUN and uses the ASTRAN placement
objective $\Phi$ directly as the reward, optimising physical layout quality.}
\label{fig:training}
\end{figure*}

\begin{algorithm}[t]
\caption{\method{} / \SGPN Training (Phases 1--2)}
\label{alg:train}
\begin{algorithmic}[1]
\Require Dataset $\mathcal{D}$, function list $\mathcal{F}$,
         hyperparameters $K, M, \lambda_H, \varepsilon_{\mathrm{clip}}$
\State Pretrain policy $\pi_\theta$ by NLL on $\mathcal{D}$ (Phase 1)
\State Initialize frozen agent $\pi_{\mathrm{old}} \leftarrow \pi_\theta$
\For{each epoch}
  \State Sample batch of functions $\mathcal{B} \subseteq \mathcal{F}$
  \For{each $f \in \mathcal{B}$}
    \State Generate $K$ trajectories using $\pi_{\mathrm{old}}$
    \State Update $T^*(f)$ if shorter complete circuit found
    \State Compute rewards $\{r_k\}$ and GRPO advantages $\{\hat{A}_k\}$
  \EndFor
  \State Compute batched PPO + entropy loss
  \State Update $\theta$ via Adam
  \If{step $\equiv 0 \pmod{M}$} $\pi_{\mathrm{old}} \leftarrow \pi_\theta$ \EndIf
\EndFor
\end{algorithmic}
\end{algorithm}

\subsection{Network Reduction}
\label{sec:reduction}

A generated network may contain transistors that do not affect the function
it computes, and counting them overstates its cost.  We therefore reduce
every network before scoring it, using two nested notions.
A transistor is \emph{structurally dead} if it lies on no simple
\textsc{src}$\to$\textsc{snk} path; such devices can never carry current and
are removed by keeping only those edges whose biconnected block lies on the
block-cut-tree path between the rails.  More generally, a network is
\emph{irredundant} if no single transistor can be deleted while still
covering every on-pattern; we delete such transistors greedily to a fixed
point.

\begin{lemma}[Deletion preserves safety]
\label{lem:deletion}
Let $G' \subseteq G$ be obtained by deleting transistors.  If $G$ conducts
under no off-pattern, then neither does $G'$.
\end{lemma}
Lemma~\ref{lem:deletion} is what makes reduction cheap: deletion can only
\emph{reduce} conduction, so no off-pattern re-verification is needed and it
suffices to re-check on-pattern coverage.  We write $T_{\mathrm{eff}}$ for
the reduced size and use it both when reporting results and inside the
reward, so the policy is never penalized for redundancy that a downstream
pass would remove.  This matters quantitatively: on the $6$-input NSP benchmark
of Section~\ref{sec:exp}, reduction lowers mean generated size from $41.4$
to $19.0$ transistors (worked example in App.~A).

\subsection{Phase 3: Physical-Aware Joint PDN+PUN Synthesis}
\label{sec:physical}

Phase~2 optimizes transistor count for a single switching network.
Phase~3 extends the RL loop to \emph{joint} synthesis of the pull-down
network (PDN) and pull-up network (PUN) of a complete CMOS cell,
with a physical placement reward from ASTRAN~\cite{astran}.

\paragraph{PDN and PUN function derivation.}
For a CMOS cell implementing $f(\mathbf{x})$, both networks must be
synthesized simultaneously:
\begin{itemize}[leftmargin=*, topsep=1pt, itemsep=0pt]
  \item \textbf{PDN} (NMOS, pulls output to GND when $f=0$):
        switching function $= \lnot f(\mathbf{x})$;
        on-patterns $= \Poff$ of $f$.
  \item \textbf{PUN} (PMOS, pulls output to VDD when $f=1$):
        switching function $= f(\lnot\mathbf{x})$, because PMOS
        transistors conduct when gate $= 0$;
        on-patterns $= \{\lnot\pi : \pi \in \Pon\}$.
\end{itemize}
Both are presented to the \emph{same} \method{} policy as distinct
switching-network synthesis problems: the policy is conditioned
on on/off patterns and is agnostic to transistor type (NMOS/PMOS).

\paragraph{SPICE generation.}
A complete PDN+PUN pair is encoded as a standard SPICE \texttt{.subckt}:
NMOS transistors (\texttt{NMOS\_VTL}, $W=90\,$nm, $L=50\,$nm) implement
the PDN with source at \texttt{GND}, drain at the output \texttt{ZN},
and internal nodes \texttt{nd$k$}; PMOS transistors implement the PUN
symmetrically.
Gate signals map directly from literal labels: positive literal $x_i$
$\to$ port \texttt{$x_i$}; negative literal $\lnot x_i$
$\to$ complemented port \texttt{$x_i$\_N}.
Power nets are named \texttt{VCC}/\texttt{GND} to match ASTRAN's
default supply identifiers.

\paragraph{Physical reward via ASTRAN.}
ASTRAN places the combined SPICE netlist using Threshold Accepting
(parameters: fold 2~0, SA quality~1, 1~attempt).
It returns five placement metrics on the ``Final cost'' line:
cell width $W_{\mathrm{cpp}}$ (CPP = total poly columns),
gate mismatches $N_{\mathrm{gm}}$, wire length $W_{\mathrm{wl}}$,
routing density $D_{\mathrm{rt}}$, and diffusion cuts $N_{\mathrm{gap}}$.
We define the physical objective as ASTRAN's own weighted sum:
\begin{equation}
  \Phi = w_c W_{\mathrm{cpp}} + w_m N_{\mathrm{gm}}
       + w_w W_{\mathrm{wl}} + w_d D_{\mathrm{rt}}
       + w_g N_{\mathrm{gap}},
  \label{eq:phi}
\end{equation}
with weights $(w_c, w_m, w_w, w_d, w_g) = (3, 4, 1, 4, 2)$ matching
the ASTRAN SA cost function exactly, so minimizing $\Phi$
is equivalent to minimizing ASTRAN's internal objective.

The joint reward for a complete PDN+PUN pair is:
\begin{equation}
  r = \frac{\Phi^*(f)}{\Phi}, \qquad
  \Phi^*(f) = \min_{\text{all pairs seen for }f} \Phi,
  \label{eq:physical_reward}
\end{equation}
where $\Phi^*(f)$ is the self-tightening best objective found so far.
Incomplete pairs (either PDN or PUN fails to cover all on-patterns
within $T_{\max}$ steps) receive $r=0$.

\paragraph{Paired trajectory collection.}
For each function $f$ per PPO step, $K$ paired trajectories
$({\tau_k^{\mathrm{PDN}}, \tau_k^{\mathrm{PUN}}})_{k=1}^K$ are generated.
ASTRAN is called only for complete pairs (up to $K$ calls per function per
step), and calls are parallelized across a thread pool.
The GRPO advantage normalizes rewards within the $K$ pairs:
$\hat{A}_k = (r_k - \bar{r}) / (\sigma_r + \varepsilon)$.
The same scalar $\hat{A}_k$ is applied to every step in both
$\tau_k^{\mathrm{PDN}}$ and $\tau_k^{\mathrm{PUN}}$:
if the pair achieves better-than-average placement, both networks
are reinforced; otherwise both are penalized.

\begin{algorithm}[t]
\caption{\method{} Phase~3: Physical-Aware Joint PDN+PUN RL}
\label{alg:physical}
\begin{algorithmic}[1]
\Require Phase~2 policy $\pi_\theta$, function list $\mathcal{F}$,
         ASTRAN binary, weights $(w_c, w_m, w_w, w_d, w_g)$
\State Initialize $\Phi^*(f) \leftarrow \infty$ for all $f$
\State $\pi_{\mathrm{old}} \leftarrow \pi_\theta$
\For{each epoch}
  \State Sample batch $\mathcal{B} \subseteq \mathcal{F}$
  \For{each $f \in \mathcal{B}$}
    \For{$k = 1, \ldots, K$}
      \State $\tau_k^{\mathrm{PDN}} \leftarrow$ generate($\pi_{\mathrm{old}}$, $\lnot f(\mathbf{x})$)
      \State $\tau_k^{\mathrm{PUN}} \leftarrow$ generate($\pi_{\mathrm{old}}$, $f(\lnot\mathbf{x})$)
      \If{both complete}
        \State $\Phi_k \leftarrow$ ASTRAN(SPICE($\tau_k^{\mathrm{PDN}}, \tau_k^{\mathrm{PUN}}$))
        \State Update $\Phi^*(f)$ if $\Phi_k < \Phi^*(f)$
        \State $r_k \leftarrow \Phi^*(f) / \Phi_k$
      \Else\ $r_k \leftarrow 0$
      \EndIf
    \EndFor
    \State GRPO: $\hat{A}_k \leftarrow (r_k - \bar{r}) / (\sigma_r + \varepsilon)$
  \EndFor
  \State Update $\theta$ via PPO + entropy loss on all steps
  \If{step $\equiv 0 \pmod{M}$} $\pi_{\mathrm{old}} \leftarrow \pi_\theta$ \EndIf
\EndFor
\end{algorithmic}
\end{algorithm}
